\chapter{Обход графа: один движок, три политики}
\label{ch:walker}

\section{Случайное блуждание как универсальный механизм}

Почти всё, что Рин делает со своей памятью, --- это обход графа роем случайных
блуждающих агентов, распространяющих активацию. Retrieval, внутренняя жизнь
и сон --- не три разных движка, а три политики одного движка, различающиеся
смещением перехода, вероятностью телепорта и набором операций.

Общее правило перехода --- нормированное произведение факторов:
\begin{equation}
P(A \to B) \propto
w_{AB}
\cdot \mathrm{rel}(B \mid q)
\cdot \mathrm{mc}(B \mid \mood)
\cdot \mathrm{unc}(B)^{\beta_u},
\label{eq:trans}
\end{equation}
где $\mathrm{mc}$ --- из \eqref{eq:moodcong}, а показатель $\beta_u$
включает/выключает любопытство.

\begin{intuition}[Что такое «температура»]
«Температура» --- одна ручка: насколько выбор следующего шага случаен против
жадного. Формально переход берётся softmax-ом от логарифма правой части
\eqref{eq:trans}, делённого на температуру $T$: при $T \to 0$ --- жадный шаг по
сильнейшему ребру (эксплуатация), при большом $T$ --- почти равновероятный
выбор соседа (исследование). Термин из имитации отжига; мистики нет.
\end{intuition}

\section{Три политики как параметризация}

\begin{center}
\begin{tabular}{>{\raggedright\arraybackslash}p{0.30\textwidth} c c c >{\raggedright\arraybackslash}p{0.25\textwidth}}
\toprule
Политика & $\alpha$ (телепорт) & $\beta_u$ & $T$ & Что запускает \\
\midrule
retrieval (цель) & $0$ & $0$ & низкая & извлечение под запрос \\
внутренняя жизнь (любопытство) & мал. & $>0$ & средняя & inspect, fusion у неопределённых \\
сон (покрытие) & большой & $0$ & высокая & fusion/inspect у случайных \\
\bottomrule
\end{tabular}
\end{center}

\smallskip\noindent\textbf{Retrieval --- целенаправленная.}
Стартует от заякоренных запросом узлов, собирает окрестность активированного
содержания. Подчеркнём (глава~\ref{ch:personality}): обход собирает ингредиенты,
а не готовую реакцию.

\smallskip\noindent\textbf{Внутренняя жизнь --- по неопределённости.}
Смещена к узлам с высоким \texttt{need\_to\_discover}, низкой уверенностью,
новизной ($\beta_u > 0$). «Любопытство» --- целенаправленный поиск
неопределённости, не случайность.

\smallskip\noindent\textbf{Сон --- телепортирующий обходчик.}
Политика стохастического покрытия: с вероятностью $\alpha$ телепорт в
случайный узел, иначе --- шаг по \eqref{eq:trans}. Это случайное блуждание с
рестартом (персонализированный PageRank). Его стационарное распределение
удовлетворяет
\begin{equation}
\pi = (1 - \alpha) P^{\top}\pi + \alpha r,
\label{eq:ppr}
\end{equation}
где $r$ --- распределение рестарта. Телепорт делает блуждание эргодическим ---
это ключевое свойство для эпистемической гигиены (глава~\ref{ch:topology}).

\begin{intuition}[Три способа гулять по городу]
Retrieval --- идти по адресу (есть цель, выбираешь осмысленный путь).
Внутренняя жизнь --- идти туда, где интересно и непонятно (тянет к незнакомым
переулкам). Сон --- время от времени садиться в случайное такси и ехать в
случайный район: не потому, что там дело, а чтобы не застревать в одном
квартале. Один пешеход, три правила выбора улицы.
\end{intuition}

\section{Что walker делает с графом: читает, предлагает, но не коммитит}
\label{sec:proposecommit}

Естественный вопрос: обходчик сам правит граф --- сливает узлы, меняет
уверенности, ретрактит? Ответ, определяющий всю конкуренцию в системе: нет.
Walker --- это сенсор-предлагатель, а не писатель. Мы разводим три роли,
которые легко слить в одну и потом мучиться с гонками:
\begin{itemize}
\item чтение (обход) --- walker ходит по графу и собирает активированное
содержание; это то, чем занимаются все три политики;
\item предложение --- итог обхода не «правка», а типизированное предложение
операции: не TODO-текст, а структура вида
\texttt{\{op: consolidate, members:[A, B, C, D], rationale, support\_provenance, walker\_id, t\}},
которая кладётся в буфер предложений;
\item коммит --- отдельная, сериализованная стадия (тот же
schema-writer / оператор консолидации, гл.~\ref{ch:fusion}) вычитывает буфер и
применяет или отклоняет операции, дописывая событие в лог
(гл.~\ref{ch:runtime}, раздел~\ref{sec:eventlog}).
\end{itemize}

Это ровно наш принцип «выход fusion --- гипотеза, а не факт»,
распространённый на все структурные операции: обход рождает кандидата,
коммиттер его судит. Зачем такая дисциплина, а не «walker сам чинит граф на
месте»: во-первых, конкурентная безопасность --- структурный писатель ровно
один, значит двух одновременных структурных правок не бывает по построению
(гонки, гл.~\ref{ch:runtime}); во-вторых, аудируемость --- каждая правка проходит
через event-log и объяснима реплеем; в-третьих, эпистемика --- предложение это
гипотеза, а коммит это отдельный гейтируемый шаг, где работает независимость
свидетельств и карантин (гл.~\ref{ch:epistemics}). Walker предлагает ---
коммиттер располагает.

\begin{intuition}[Единственное, что walker пишет сам --- но это тоже запись]
Полный запрет на запись был бы неудобен: обход должен оставлять след «здесь
были, это подсветили». Поэтому walker напрямую трогает только поле активации /
яркости $a_v$ (и связанную телеметрию внимания). Точная формулировка важна во
избежание тихих багов: это не «не запись», а неструктурная коммутативная
запись --- не узлы, не рёбра, не уверенности, а состояние внимания, которому
всё равно нужны TTL, затухание и \texttt{access\_reason} (гл.~\ref{ch:timescales}).
«Коммутативная» значит, что итог не зависит от порядка параллельных записей.
Но конкретно для яркости не годится аддитивный счётчик: на хабах, через которые
проходят сотни walker'ов, $a_v \mathrel{+}= 1$ улетает в космос. Берём max-CRDT:
$a_v \leftarrow \max(a_v, e_{rw})$, где $e_{rw}$ --- энергия зашедшего обходчика.
Тогда яркость ограничена сверху энергией агента, спад предсказуем, а сотня
одновременных подсветок даёт корректный итог без блокировок. Всё, что
структурно, --- только через предложение и коммит.
\end{intuition}

\section{Фоновая разведка: один процесс, а не «день/ночь»}
\label{sec:background}

Разделение на «дневное любопытство» и «ночной обходчик» --- удобная, но
неточная рамка. Точнее: есть один всегда включённый фоновый разведчик, а
«день» и «ночь» --- два конца континуума по свободным ресурсам, а не
переключатель.

Формализуем «мощности, внимание, фокус». Пусть $\rho(t) \in [0, 1]$ ---
foreground-нагрузка: насколько текущий момент требует активного ответа
($\rho \approx 1$ в напряжённом диалоге, $\rho \approx 0$ в простое). Общий
бюджет $B_{\text{total}}$ делится с приоритетом переднего плана; фоновой
разведке достаётся
\[
B_{\text{bg}} = B_{\text{total}} (1 - \rho).
\]

Фокус --- это активное множество $A_{\text{focus}}$ (узлы высокой активации
$a_v$). Параметры фонового роя --- функции $B_{\text{bg}}$: число walker'ов
$\propto B_{\text{bg}}$; доля телепортов $\alpha$ и «широта» растут при падении
$\rho$; стартовое распределение интерполирует
\[
r = (1 - \beta)\cdot \mathrm{anchor}(A_{\text{focus}}) + \beta\cdot \mathrm{uniform},
\qquad \beta \uparrow \text{ при } \rho \downarrow.
\]

Тогда прежние «любопытство» и «сон» --- это два режима одного процесса:
высокий $\rho$ (мало бюджета) $\Rightarrow$ мало walker'ов, $\alpha \approx 0$,
старт от фокуса --- локальная разведка вокруг текущей темы; низкий $\rho$
(много бюджета) $\Rightarrow$ много walker'ов, большой $\alpha$, старт почти
равномерный --- широкий телепорт-скраб (глава~\ref{ch:sleep}).

\begin{intuition}[Мысль на перекуре]
Именно поэтому яркая ассоциация может прийти днём, «на ровном месте»:
фоновый разведчик работает всегда, и даже на скудном дневном бюджете иногда
набредает на далёкую связь. Развёрнутый пример такой дневной находки ---
в главе~\ref{ch:fusion} (House / Cast Away).
\end{intuition}

\subsection{Поля активации: быстрое внимание и медленный след --- не путать}
\label{sec:brightness}

Здесь легко совершить ошибку --- свалить в одно поле «яркость», «свежесть»,
«долговременную силу» и «важность». Соблазнительно, но потом это поле начнёт
делать слишком много, и, что хуже, диагностика начнёт тренировать память:
скраб, дебаг и тесты будут поднимать ту же переменную, что и настоящее
вспоминание, и случайно укреплять узлы. Поэтому разведём разные по темпу
состояния:
\begin{itemize}
\item $a_{\text{work}}(v,t) \in [0,1]$ --- рабочая яркость внимания. Живёт
секунды-минуты, формирует зону фокуса $A_{\text{focus}}$ и вход дневных
walker'ов, быстро гаснет. Это неструктурная коммутативная запись обходчика
(max-CRDT выше).
\item $B_{\text{base}}(v)$ --- медленная сила следа (это и есть ACT-R base-level,
гл.~\ref{ch:timescales}). Живёт дни-месяцы, обновляется только настоящими
значимыми обращениями (по \texttt{access\_reason}), влияет на decay / забывание.
\item $\mathrm{importance}(v)$ --- значимость от appraisal / домена / структуры
(гл.~\ref{ch:personality}), не путать с частотой обращений.
\item $\mathrm{coverage\_deficit}(v)$ --- «давно не инспектировали», рычаг
ночного скраба.
\end{itemize}

Ключевое разделение: настоящее вспоминание поднимает и $a_{\text{work}}$
(мгновенно), и $B_{\text{base}}$ (чуть-чуть, медленно); а скраб / дебаг / тест
трогают в лучшем случае $\mathrm{coverage\_deficit}$ и никогда не укрепляют
след. Правило быстрой яркости:
\[
a_{\text{work}}(v,t+\Delta t) = a_{\text{work}}(v,t)e^{-\Delta t/\tau(v)}
\quad \text{(между визитами),}
\]
\[
a_{\text{work}}(v) \leftarrow \max\!\bigl(a_{\text{work}}(v), e_{rw}\bigr)
\quad \text{(genuine-визит),}
\]
где $\tau(v)$ не универсальна (важные / консолидированные остывают медленнее),
а $e_{rw}$ --- энергия обходчика (max, не сумма, --- иначе хабы взрываются).

\begin{openq}[max теряет «навязчивость» --- возможно, нужен второй канал]
Честный компромисс max: он ограничивает яркость (хорошо против хабов), но
теряет повторный фокус в коротком окне. Сто обходчиков через хаб --- max
правильно не взрывает; но если один узел многократно всплывает в текущем
рассуждении, иногда важно знать не только пик, но и интегральную экспозицию
(«это навязчиво лезет прямо сейчас»). Возможное решение --- два канала:
$a^{\text{peak}}_{\text{work}} = \max(\ldots)$ для яркости и отдельный
\texttt{visit\_mass\_window} для краткосрочной статистики (он не укрепляет
$B_{\text{base}}$ автоматически). На Stage 1 не обязателен; помечаем как
возможное расширение.
\end{openq}

Теперь смещение входа обходчиков. Распределение рестарта $r$ из
\eqref{eq:ppr} берётся не равномерным, и день с ночью различаются формой
смещения, а не разными движками:
\begin{itemize}
\item День (высокий $\rho$, радиус $r_0$): $r \propto a_{\text{work}} \times
\mathrm{relevance} \times \mathrm{uncertainty}$. Обходчик эксплуатирует текущий
фокус --- близкие ассоциации «на сейчас», далеко не отходит.
\item Ночь (низкий $\rho$, радиус $r_1 > r_0$): $r \propto \mathrm{uniform} +
\mathrm{coverage\_deficit} + \mathrm{need\_to\_revise} + \text{low-}\pi$ регионы.
Именно так, а не просто «шире от дневного фокуса» --- иначе сон останется
заложником дневного внимания. Сон нужен, чтобы уйти туда, куда внимание
никогда не ходит: anti-entropy scrub, телепорт-эргодичность, проверка
противоречий (гл.~\ref{ch:sleep}).
\end{itemize}

Лаконично: день --- эксплойт локально у горячего, ночь --- исследуй заброшенное.
Поля разные, но связанные: $a_{\text{work}}$ ведёт дневной фокус,
$\mathrm{coverage\_deficit}$ и низкая $\pi$ ведут ночь, $B_{\text{base}}$ решает,
что переживёт забывание.

\begin{intuition}[Мотыльки на остывающий свет: откуда дневные догадки]
Если дневной ретрив --- «слепой матричный ток» (раздел~\ref{sec:semgate}),
откуда берутся внезапные дневные озарения? Не из matvec. После дневного ответа
подсвеченные узлы $a_{\text{work}}$ гаснут экспоненциально минутами-часами ---
остаётся тёплое послесвечение. Одиночные фоновые RW (тот самый остаточный
бюджет $B_{\text{bg}}$) через формулу рестарта
$r \propto a_{\text{work}} \times \mathrm{relevance} \times \mathrm{uncertainty}$
получают повышенную вероятность телепортироваться именно в эти недавно-горячие
зоны --- летят на остывающий свет, как мотыльки. Кружа в тёплой области, они
сталкивают концепты на своих траекториях и роняют кандидатов в thought-inbox
(гл.~\ref{ch:runtime}). Важно: роняют именно кандидатов (в основном шум),
которых обязан просеять salience-гейт ($|\Delta \mood|$, ошибка предсказания) ---
иначе Рин завалит себя мусорным fusion. Так разделение труда честное: matvec
подсвечивает, мотыльки ползают по послесвечению, salience решает, что вообще
достойно всплыть догадкой.
\end{intuition}

\begin{intuition}[Где «живут» обходчики]
Walker'ы не объекты, которые где-то постоянно «живут». Это фоновый процесс
планировщика с бюджетом $B_{\text{bg}}$: каждый обход рождается из
распределения рестарта $r$, проходит, оставляет предложения и подсветку --- и
гаснет. Persist между ними не сами walker'ы, а поле яркости $a_{\text{work}}$ и
буфер предложений. Спросить «где живут обходчики» --- всё равно что спросить,
где живут волны: волн как вещей нет, есть среда (поле $a_{\text{work}}$) и то,
что её возмущает (планировщик).
\end{intuition}

\begin{openq}[Как измерять foreground-нагрузку $\rho$]
$\rho$ управляет тем, насколько агрессивно разведчик забирает ресурсы, но как
её мерить в рантайме --- не зафиксировано. Кандидаты: наличие активного
пользовательского хода (есть незакрытая реплика $\Rightarrow \rho$ высока);
длина очереди foreground-задач; «важность» текущего события (по appraisal,
глава~\ref{ch:personality}); загрузка вычислителя. Вероятно, комбинация, но
веса и динамика (как быстро $\rho$ спадает в простое) пока открыты.
\end{openq}

\section{Точки входа: как query превращается в якоря}
\label{sec:entrypoints}

Ретрив у нас двухшаговый: шаг 1 --- найти точки входа (якоря) в граф, шаг 2 ---
растечься из них диффузией (раздел~\ref{sec:semgate}). Про шаг 2 сказано много,
а про шаг 1 --- только «быстрым индексом». Распишем его, потому что это первая
строчка кода ретрива.

\smallskip\noindent\textbf{Сразу уточнение терминологии.}
Query не «разбивают на эмбеддинги». Его разбивают на семантические единицы,
и уже каждую единицу отдельно эмбеддят. Это разные вещи.

\smallskip\noindent\textbf{Кто разбивает --- не LLM.}
Ретрив --- рефлекс (System 1, ниже LLM-слота, гл.~\ref{ch:runtime}), поэтому
превращение query в якоря должно быть дешёвым и LLM-free на горячем пути.
Делает это тот же экстрактор, что и на write-path
(\texttt{raw\_event -> parse/extract candidates}, раздел~\ref{sec:writepath}),
только без коммита. Красивая симметрия: понимание запроса = извлечение
кандидатов из события минус запись. На Stage 1 экстрактор простой: NER +
ключевые фразы + лемматизация (для русского морфология обязательна ---
«плавки / плавок / плавками» сводятся к лемме). LLM подключается к
переформулировке мутного запроса только на эскалации --- редкий System 2, а не
норма.

\smallskip\noindent\textbf{Два сорта якорей.}
Экстрактор выдаёт две разновидности, и ищутся они по-разному:
\begin{itemize}
\item Сущностные / лексические --- именованные сущности и конкретные термины
(«Азоры», «плавки», «9 сентября»). Им нужен точный матч, а не «похожее по
вектору».
\item Холистический --- эмбеддинг всего запроса (или его клауз) для
ассоциативной диффузии.
\end{itemize}

Если единиц несколько --- это мульти-векторный ретрив, а не одно усреднённое.
Mean-pool «плавки на Азоры 9 сентября» размазал бы три резкие сущности в кашу;
поэтому каждую единицу сажаем отдельным seed'ом, а узлы, подсвеченные
несколькими seed'ами, получают буст (пересечение = релевантнее).

\smallskip\noindent\textbf{Чем эмбеддим.}
Тем же энкодером, что и \texttt{node.embedding} при записи (гл.~\ref{ch:substrate}).
Это не опция: косинус осмыслен, только если query и узлы в одном пространстве.
Вектор узла уже кэширован (считался один раз коммиттером), так что на шаге
ретрива через энкодер гоним только query-единицы, узлы не трогаем.

\begin{flushleft}\ttfamily
1 units      = extract(query + active\_focus)      \# NER + keyphrases + леммы (+ кореф)\\
2 anchors    = \{\}\\
3 for u in units:\\
4 \quad if u.is\_named\_entity:\\
5 \qquad anchors |= exact\_index.lookup(lemma(u)) \# точный/alias-индекс (forensic-подобный)\\
6 \quad anchors |= ann\_index.knn(embed(u), k)     \# вектор-ANN (HNSW), top-k по косинусу\\
7 anchors    = resolve\_entities(anchors)          \# упоминания -> канонические узлы\\
8 r          = weight\_by\_match\_confidence(anchors) \# точное > сильный вектор > слабый\\
9 \# r -> распределение рестарта для локальной разреженной диффузии (шаг 2)
\end{flushleft}

То есть точки входа приходят из двух индексов: точного / alias (для именованных
сущностей --- «Азоры» должны попасть в тот самый узел) и вектор-ANN (для
семантической близости). Их объединение, взвешенное по уверенности матча, и
есть seed-набор. Плюс entity resolution: разные упоминания сводятся к
каноническому узлу, чтобы «плавки» били в один узел, а не в пять эпизодов
покупки.

\smallskip\noindent\textbf{Дальше.}
Seed-набор становится начальным впрыском $a(0)=s(0)$ для локальной затухающей
диффузии (fast-path вокруг якорей, раздел~\ref{sec:semgate}), а расширение
гейтируется условным throttle. Важно: retrieval --- это политика с телепортом
$\alpha = 0$ (таблица политик выше), то есть источник впрыскивается один раз и
импульс затухает; рестарт (член $s(t>0)$, он же $\alpha>0$ в \eqref{eq:ppr}) ---
это сон, а не retrieval. Цепочка целиком: экстрактор дал якоря $\to$
ANN+точный индекс нашли точки входа $\to$ импульс затух $\to$ packer срезал по
порогу.

\begin{openq}[Что на Stage 1 сознательно упрощаем]
Открыты: гранулярность разбиения (по сущностям / клаузам / предложениям ---
тюнить на бенче); кореференс --- «поищи их» требует обогатить query текущим
фокусом $a_{\text{work}}$ до эмбеддинга, иначе якорей нет; порог эскалации к
LLM-переформулировке; качество alias-резолвинга под русскую морфологию.
Механика (два сорта якорей, два индекса, мульти-seed, resolution) ---
зафиксирована.
\end{openq}

\section{Семантический вентиль: дёшево и без протечек по омонимии}
\label{sec:semgate}

В \eqref{eq:trans} множитель $\mathrm{rel}(B \mid q)$ --- это семантический
вентиль: он не даёт фронту расползтись из зоны запроса. Возникает законный
страх за производительность: если на каждом шаге RW считать сходство контента
узла с $q$ заново, это катастрофа. Ответ: на лету ничего не эмбеддим.
Эмбеддинг узла считается один раз при записи (коммиттером) и лежит в поле
\texttt{embedding} (гл.~\ref{ch:substrate}). Тогда вентиль на шаге --- это
скалярное произведение фиксированного вектора $q$ на кэшированный вектор узла,
$O(d)$; а раз ретрив у нас --- усечённая локальная затухающая диффузия из
seed'ов (политика retrieval: телепорт $\alpha = 0$, источник впрыскивается один
раз --- это не PPR; рестарт-форма \eqref{eq:ppr} зарезервирована за сном), весь
фронт считается одним matvec, где вектор активации на каждом шаге умножается на
маску сходства с $q$. Как только значения падают ниже порога $\varepsilon$ ---
граница замкнулась. Никаких рекурсивно плодящихся агентов: «группа вспышек,
слабеющих к окраинам» --- это результат диффузии, а не объектная механика.

Инженерная честность про этот дневной fast-path: matvec идёт не по всему
графу. Умножать $\tilde{P}^{\top}$ на $|V|$-мерный вектор на графе в сотни
тысяч узлов --- лечь. Это локальная разреженная диффузия на ограниченном
носителе: 2--3 хопа вокруг seed-якорей, найденных ANN-индексом (HNSW), ---
субматрица смежности окрестности, а не всё $V$-пространство. И тайминги честные
--- миллисекунды, не микросекунды (основное время --- ANN-поиск якорей и
context-packer), а настоящее бутылочное горло рантайма всё равно LLM-инференс
(секунды), гл.~\ref{ch:runtime}.

Две поправки, чтобы вентиль не навредил. Первая --- жёсткость по типу
ретрива: для ordinary-ретрива режем жёстко (защита от мусора), для
creative / curiosity --- мягко (именно далёкие странные мосты нужны для fusion),
для forensic --- вентиля нет вообще, это не walker, а точный индекс-lookup
(таксономия ретрива, гл.~\ref{ch:runtime}). Слишком жёсткий вентиль убил бы
серендипность. Вторая --- полисемия: граница возникает не только из низкого
сходства с $q$, но и из sense/type-несовпадения. Классический пример
«кружка $\to$ ручка $\to$ дверь» не должен решаться тем, что «дверь слабо
похожа на кружку»: раньше ломается слой смысла. «Ручка кружки», «дверная
ручка», «ручка-перо» --- это разные sense/concept-узлы (или хотя бы
квалифицированные рёбра, гл.~\ref{ch:embryonic}). Без этого ретрив постоянно
течёт по омонимии --- а для русского это особенно больно. Значит вентиль
закрывается по $\max(\text{low-sim}, \text{sense/type-mismatch})$, а не по одной
similarity.

\subsection{Вентиль --- условный throttle, а не плоский множитель}
\label{sec:semgate-conditional}

Есть и третья поправка, и она важнейшая: как написано выше,
$\mathrm{rel}(B \mid q)$ входит в \eqref{eq:trans} множителем --- а значит
зануляет путь при низком сходстве. Это тихо убивает личность. Пример: триггер
«гирлянда». Узлы [свет], [провода] семантически близки и проходят; а узел
личной памяти [окно в Тбилиси] имеет сходство $\approx 0$ --- и множитель его
гасит, Рин выдаёт энциклопедию. В живой голове [Тбилиси] вспыхивает первым:
у личного ребра колоссальный вес / валентность, несмотря на семантическую
даль. Значит вентиль должен быть условным throttle, а не вето:
\[
\mathrm{Score}(A \to B) =
f(w_{AB}, \mathrm{valence}_B)\cdot
\mathrm{Throttle}(B,q)\cdot
\mathrm{mc}(B \mid \mood)\cdot
\mathrm{unc}(B)^{\beta_u},
\]
где
\[
\mathrm{Throttle}(B,q) =
\begin{cases}
1, & \shortstack[l]{если ребро из личного слоя (session/life)$;$\\
или $|\mathrm{valence}_B| > \theta_{\text{emotion}}$,}\\
\mathrm{rel}(B \mid q), & \text{иначе (обычное знание --- тормоз против дрейфа).}
\end{cases}
\]

То есть сильное личное или аффективно заряженное ребро --- «синаптический
хайвей» --- проходится независимо от сходства с запросом; а throttle тормозит
только слабые / родовые рёбра, чтобы не было случайного топологического
дрейфа.

Три предохранителя, без которых это сломается. Не по сырому $w_{AB}$:
пробивать вентиль порогом веса нельзя --- у эмбриональных рёбер (ConceptNet)
веса из коробки огромны ($w \approx 2{,}5$), и Рин уплыла бы по банальной
таксономии. Поэтому гейт по происхождению ребра (личный слой, не врождённый
каркас), а не по величине веса. $f(w_{AB}, \mathrm{valence})$ всё равно
ранжирует: $\mathrm{Throttle}=1$ не значит «все личные узлы хлынут» ---
Тбилиси всплывает не потому, что ребро личное, а потому, что оно личное и с
огромным $w_{AB}$; слабое личное ребро с $\mathrm{Throttle}=1$ всё равно
проиграет по $f$. Стоп-условие три-гейта: путь глохнет, только если
одновременно низкое сходство и слабое ребро и нейтральная валентность; сильна
любая из трёх --- агент проходит. Сухой-но-важный boundary («не шути про
доту») проходит через layer-гейт даже при низкой валентности. Омонимию
(кружка $\to$ ручка $\to$ дверь) по-прежнему режет отдельный sense/type-гейт,
а не throttle.

\begin{intuition}[Прайминг: подсвечено, но не сказано]
База при этом не пропадает: «у гирлянды есть провода» матричная диффузия
подсвечивает (ток туда течёт) --- это прайминг. Но у этих узлов почти ноль
ситуативной значимости и новизны, поэтому context-packer их сознательно не
выносит наверх (раздел~\ref{sec:packer}). Они «тёплые» на фоне сознания, но не
доходят до языка --- ровно как человеческое подсознание держит очевидное
наготове, не проговаривая.
\end{intuition}

\begin{principle}[Вентиль на эмбеддингах --- фильтр родственности, не релевантности]
Эмпирически выяснено (стенд, \texttt{reports/experiment\_3--4}): множитель
$\mathrm{rel}(B \mid q)$, построенный на дистрибутивных эмбеддингах (косинус к
вектору запроса), фильтрует родственность (co-occurrence), а не причинную
релевантность. Он давит off-topic дрейф (годится для запросов вида «что есть
$X$» / «кто из этих $X$»), но на «почему»-запросах систематически
недобирает: узел может быть ближе к слову запроса как ассоциат, не будучи
ответом («beverage» ближе к «energy», чем «caffeine», из-за маркетингового
со-упоминания). Это ограничение инструмента, а не бага; для причинного выбора
нужен не косинус, а рассуждение.
\end{principle}

\begin{principle}[Конвергенция на хабе пересиливает per-edge релевантность; лечится редкостью, не гейтом]
Эмпирически выяснено (там же): при диффузии узел, достижимый из нескольких
якорей сразу, накапливает их вклады, и это суммирование может перевесить более
высокий per-edge $\mathrm{rel}(B \mid q)$ у односвязного правильного ответа ---
даже когда вентиль оценил ответ выше. То есть провал «генеричный дистрактор
победил» решает конвергенция на хабе, а не релевантность; это тот же зверь,
что концентрация $\pi$ на хабах (гл.~\ref{ch:topology}). Значит корректив ---
отдельный штраф за редкость / степень (IDF-подобный: генеричный общий сосед
давится), а не докрутка вентиля. Один такой член чинит сразу три наблюдения:
дистрактор-хаб, таксономический дрейф ретрива и хаб-концентрацию.
\end{principle}

\begin{principle}[Ретрив отвечает за recall, LLM --- за причинный выбор]
Отсюда критерий успеха ретрива: не «правильный узел первый в ранге», а
«правильный узел попал в упакованный контекст в рамках бюджета при приемлемой
доле дистракторов». Граф --- ассоциативная память (поднять кандидата в
окрестность); причинное ранжирование («из этих caffeine --- и есть почему»)
делает LLM на собранном контексте. Требовать от графа самому назвать причинный
ответ --- значит возлагать рассуждение на recall.
\end{principle}

\section{Единство и следствие}

Один субстрат, один блуждающий агент, три параметризации: смещение перехода,
вероятность телепорта, набор операций. Реализуется одна машина обхода с
параметрами (псевдокод --- приложение~\ref{ch:algorithms}). Родственные
механизмы: spreading activation (Collins--Loftus) как модель ассоциативного
припоминания и персонализированный PageRank / RWR как его вычислительная
форма; в волте RIN --- под именами SYNAPSE и GAAMA.
